% Part: sets-functions-relations
% Chapter: sets
% Section: russells-paradox

\documentclass[../../../include/open-logic-section]{subfiles}


\begin{document}

\olfileid{sfr}{set}{rus}
\olsection{રસેલનો વિરોધાભાસ}

ઘટકો દ્વારા નિર્ધારિત સમાનતાનો સિદ્ધાંત $\phi(x)$ હોય તેવા $x$ના
\emph{ચોક્કસ એ જ} ગણ માટે $\Setabs{x}{\phi(x)}$ સંજ્ઞાને યોગ્ય ઠરાવે છે.
પરંતુ આ સિદ્ધાંત \emph{ખરેખર} માત્ર નીચેના વિચારને યોગ્ય ઠરાવે છે:
\emph{જો} એવો ગણ હોય જેના સભ્યો બરાબર $\phi$ ગુણધર્મ ધરાવતી બધી
વસ્તુઓ જ હોય, \emph{તો} આવો એક જ ગણ હોય.
બીજી રીતે કહીએ તો, કોઈ~$\phi$ નક્કી કર્યા પછી,
$\Setabs{x}{\phi(x)}$ ગણ \emph{અસ્તિત્વ ધરાવતો હોય તો} અનન્ય છે.

પરંતુ આ શરત મહત્ત્વની છે!{} ખાસ કરીને, દરેક ગુણધર્મ વડે
\emph{ગણરચના} કરી શકાતી નથી. એટલે કે કેટલાક ગુણધર્મો ગણની વ્યાખ્યા
આપતા \emph{નથી}. જો બધા ગુણધર્મો આમ કરતા હોત, તો સીધો જ
વિરોધ ઊભો થાત. આનું સૌથી જાણીતું ઉદાહરણ રસેલનો વિરોધાભાસ છે.

ગણો બીજા ગણોના !!{element}s હોઈ શકે છે---દાખલા તરીકે,
ગણ~$A$નો ઘાતગણ ગણોથી બનેલો છે.
આથી કોઈ ગણ બીજા ગણનો !!a{element} છે કે નહીં તે પૂછવું કે તપાસવું અર્થપૂર્ણ છે.
શું કોઈ ગણ પોતાનો જ સભ્ય હોઈ શકે?
ગણના ખ્યાલમાં આને બાકાત રાખે એવું કંઈ દેખાતું નથી.
દાખલા તરીકે, જો \emph{બધા} ગણો વસ્તુઓનો સંગ્રહ બનાવતા હોય,
તો કોઈને લાગે કે તેમને એક જ ગણમાં એકત્ર કરી શકાય---બધા ગણોનો ગણ.
અને તે પોતે ગણ હોવાથી બધા ગણોના ગણનો !!a{element} હોય.

પોતે પોતાનો !!a{element} ન હોવાના, એટલે કે \emph{સ્વસભ્ય ન હોવાના},
ગુણધર્મનો વિચાર કરીએ ત્યારે રસેલનો વિરોધાભાસ ઊભો થાય છે.
ધારો કે એવા બધા ગણોનો ગણ છે જે પોતે પોતાના !!a{element} નથી.
તો શું
\[
R = \Setabs{x}{x \notin x}
\]
અસ્તિત્વ ધરાવે છે? આપણે સાબિત કરી શકીએ છીએ કે તે અસ્તિત્વ ધરાવતો નથી.

\begin{thm}[રસેલનો વિરોધાભાસ]\ollabel{thm:russells-paradox}
	$R = \Setabs{x}{x \notin x}$ એવો કોઈ ગણ નથી.
\end{thm}

\begin{proof}
જો $R = \Setabs{x}{x \notin x}$ અસ્તિત્વ ધરાવતો હોય, તો
$R \in R$ ત્યારે અને ત્યારે જ થાય જ્યારે $R \notin R$ હોય,
જે પરસ્પરવિરોધ છે.
\end{proof}

\begin{tagblock}{novice}
\begin{explain}
આ સાબિતીને વધુ ધીમેથી સમજીએ. જો $R$ અસ્તિત્વ ધરાવતો હોય, તો
$R \in R$ છે કે નહીં તે પૂછવું અર્થપૂર્ણ છે.
ધારો કે ખરેખર $R \in R$ છે. હવે $R$ની વ્યાખ્યા એવા બધા ગણોના ગણ
તરીકે આપેલી છે જે પોતે પોતાના !!{element}s નથી.
આથી, જો $R \in R$ હોય, તો $R$ પોતે $R$ને વ્યાખ્યાયિત કરતો ગુણધર્મ ધરાવતો નથી.
પરંતુ આ ગુણધર્મ ધરાવતા ગણો જ $R$માં છે, તેથી $R$ એ $R$નો
!!a{element} હોઈ શકે નહીં, એટલે કે $R \notin R$.
પરંતુ $R$ એકસાથે $R$નો !!a{element} હોય અને ન પણ હોય તે શક્ય નથી.
આથી આપણને પરસ્પરવિરોધ મળે છે.

$R \in R$ની ધારણા પરસ્પરવિરોધ તરફ દોરી જાય છે, તેથી $R \notin R$ છે.
પરંતુ આ પણ પરસ્પરવિરોધ તરફ દોરી જાય છે!{}
કારણ કે જો $R \notin R$ હોય, તો $R$ પોતે $R$ને વ્યાખ્યાયિત કરતો ગુણધર્મ
ધરાવે છે, અને તેથી, બીજા બધા સ્વસભ્ય ન હોય તેવા ગણોની જેમ,
$R$ પણ $R$નો !!a{element} હોય.
ફરીથી, તે એકસાથે $R$નો !!a{element} ન હોય અને હોય પણ ખરો તે શક્ય નથી.
\end{explain}
\end{tagblock}

\begin{digress}
રસેલના વિરોધાભાસમાં ન ફસાતો ગણસિદ્ધાંત કેવી રીતે રચી શકાય?
એટલે કે $R = \Setabs{x}{x \notin x}$ અસ્તિત્વ ધરાવે છે એવો
\emph{અસંગત} દાવો ટાળતો ગણસિદ્ધાંત કેવી રીતે રચી શકાય?
તે માટે એવી પૂર્વધારણાઓ નક્કી કરવી પડે જે ગણો ક્યારે અસ્તિત્વ ધરાવે
છે (અને ક્યારે નથી ધરાવતા) તે કહેવા માટે અત્યંત ચોક્કસ શરતો આપે.

આ પ્રકરણમાં રૂપરેખા આપેલો ગણસિદ્ધાંત આવું કરતો નથી.
તે \emph{ખરેખર સહજ ખ્યાલો પર આધારિત} છે.
તે માત્ર એટલું કહે છે કે ગણો ઘટકો દ્વારા નિર્ધારિત સમાનતાના સિદ્ધાંતને
અનુસરે છે અને જો તમારી પાસે કેટલાક ગણો હોય, તો તેમનો યોગગણ,
છેદગણ વગેરે બનાવી શકો છો.
ગણસિદ્ધાંતને આના કરતાં વધુ ચોકસાઈથી વિકસાવવો શક્ય છે.
\oliflabeldef{cumul:::part}{આ ચોકસાઈ આપણે ભાગ \olref[cumul][][]{part}
માટે રાખીશું. હાલ આપણે સહજ ખ્યાલોના આધારે આગળ વધીશું
અને કાળજીપૂર્વક વિરોધો ટાળવાનો પ્રયાસ કરીશું.}{}
\end{digress}


\end{document}

